\phantomsection\addcontentsline{toc}{section}{Industry 4.0}
\ECchapter{22}{Industry 4.0}{How can connected factories improve production?}{ECGreen}

\ECObjectives{
\begin{itemize}
\item Describe the layered architecture of a connected production system.
\item Calculate and interpret availability, performance, quality and OEE.
\item Evaluate cybersecurity, traceability and human-factor requirements.
\end{itemize}}

\begin{multicols}{2}
\begin{engineeringnews}
Factories are connecting machines, robots, inspection stations and planning software. The objective is
not simply to automate more tasks, but to measure losses, adapt to variation and support faster,
better-informed decisions.
\end{engineeringnews}

\begin{ecarticle}
Industry 4.0 combines connected machines, digital models, automation and data analysis. At machine
level, sensors measure cycle time, motor current, vibration, position and product quality. Controllers
operate the process in real time, while industrial networks transmit selected information to
supervisory systems.

A connected factory is organised in layers. The physical layer contains machines and products. The
control layer contains PLCs, drives and robot controllers. A manufacturing execution system coordinates
orders, traceability and work instructions. Enterprise software manages resources, suppliers and demand.
Information must move between these layers without disturbing deterministic control or production safety.

Condition monitoring can detect a damaged bearing before failure. Vision systems can inspect every
part. Production data can reveal that a defect appears only with one tool, one material batch or one
temperature range. These applications reduce downtime and waste only when each indicator is connected
to a physical cause and a defined action.

Overall equipment effectiveness combines availability, performance and quality:
\[
\mathrm{OEE}=A\times P\times Q.
\]
A machine with high nominal speed but frequent stops may therefore have poor OEE. Engineers must locate
the dominant loss rather than improve the most visible one.

Connectivity also introduces risks. A cyberattack can modify process parameters, legacy equipment may
use insecure protocols and operators may receive too many alarms. Successful projects combine network
segmentation, clear interfaces, training and human expertise.
\end{ecarticle}

\begin{tcolorbox}[title=\sffamily\bfseries Engineering Insight,colback=yellow!10,colframe=ECOrange]
Collecting more data does not automatically improve production. A useful indicator needs a physical
meaning, a decision threshold, an owner and a verification method.
\end{tcolorbox}

\begin{engineeringfocus}
\textbf{Connected production architecture}
\begin{center}
\resizebox{.82\linewidth}{!}{%
\begin{tikzpicture}[font=\scriptsize,>=Latex,node distance=3.5mm]
\node[draw=ECGreen,fill=ECGreen!10,rounded corners] (machines) {machines and sensors};
\node[draw=ECBlue,fill=ECLightBlue,rounded corners,above=of machines] (control) {PLCs and control};
\node[draw=ECPurple,fill=ECPurple!10,rounded corners,above=of control] (mes) {MES and traceability};
\node[draw=ECOrange,fill=ECOrange!10,rounded corners,above=of mes] (erp) {planning and enterprise};
\draw[<->,thick] (machines)--(control); \draw[<->,thick] (control)--(mes); \draw[<->,thick] (mes)--(erp);
\end{tikzpicture}%
}
\end{center}
Real-time control should remain local even if an upper information layer becomes unavailable.
\end{engineeringfocus}

\begin{keyfigures}
\begin{tabularx}{\linewidth}{@{}Xr@{}}
Availability loss & downtime\\
Performance loss & slow cycles\\
Quality loss & rejected parts\\
Typical controller & PLC\\
Critical requirement & secure network
\end{tabularx}
\end{keyfigures}

\begin{tcolorbox}[title=\sffamily\bfseries OEE components,colback=white,colframe=ECGreen]
\begin{center}
\begin{tikzpicture}
\begin{axis}[ybar,bar width=8pt,width=.96\linewidth,height=47mm,ymin=0,ymax=105,ylabel={Score (\%)},symbolic x coords={Conventional,Connected},xtick=data,legend style={font=\scriptsize,at={(.5,-.22)},anchor=north,legend columns=3},ticklabel style={font=\scriptsize}]
\addplot table[x=mode,y=availability,col sep=comma]{data/EC22/oee.csv};
\addplot table[x=mode,y=performance,col sep=comma]{data/EC22/oee.csv};
\addplot table[x=mode,y=quality,col sep=comma]{data/EC22/oee.csv};
\legend{Availability,Performance,Quality}
\end{axis}
\end{tikzpicture}
\end{center}
\end{tcolorbox}

\begin{vocabulary}
\begin{tabularx}{\linewidth}{@{}lX@{}}
shop floor & atelier de production\\
PLC & automate programmable\\
downtime & temps d'arrêt\\
traceability & traçabilité\\
work instruction & instruction de travail\\
condition monitoring & surveillance d'état\\
quality control & contrôle qualité\\
throughput & cadence / débit\\
legacy machine & machine ancienne\\
cybersecurity & cybersécurité
\end{tabularx}
\end{vocabulary}

\begin{questions}
\item What are the main layers of a connected factory?
\item How can condition monitoring reduce downtime?
\item What three factors form OEE?
\item Why is collecting all possible data ineffective?
\item Why should machine control remain local?
\end{questions}

\begin{applicationexercise}
During one shift, a line has an availability of $90\%$, a performance score of $85\%$ and a quality
score of $98\%$.
\begin{enumerate}
\item Calculate the OEE.
\item Availability rises to $95\%$. Calculate the new OEE.
\item Determine the improvement in percentage points and relative percentage.
\item Decide whether improving availability or quality by one additional point produces the larger OEE gain.
\end{enumerate}
\end{applicationexercise}

\begin{engineeringchallenge}
A packaging line suffers short, frequent stops. Define a measurement plan, three indicators and their
thresholds. Draw the information path from sensor to dashboard, identify one cybersecurity risk and
specify how operators will verify that the proposed improvement actually increases OEE.
\end{engineeringchallenge}

\begin{speaking}
Will connected factories create better industrial jobs or reduce human autonomy? Discuss safety,
skills, productivity, surveillance and responsibility.
\end{speaking}
\end{multicols}

\subsection*{Sources}
\ECsource{International Society of Automation}{https://www.isa.org/}
\ECsource{NIST -- Smart Manufacturing}{https://www.nist.gov/topics/smart-manufacturing}